\section{Exact minimality in a specified reporter family}
\label{enum:section}

The twenty-bit construction is small enough to admit a finite
minimality question. This section answers that question for the
binomial-block reporters defined below. It does not establish a
minimum dimension among all observations of a Boolean cube.

\subsection{The class being classified}

\begin{definition}[Binomial sum-level reporter]\label{enum:definition}
Choose positive integers \(n,m\), and take \(m+1\) independent blocks
of \(n\) uniform signs. Write their raw sums as
\(U,Z_1,\ldots,Z_m\). Their common alphabet and probabilities are
\[
 \mathcal S_n=\{x_j=2j-n:0\le j\le n\},\qquad
 p_j=2^{-n}\binom nj.
\]
Choose tests \(I_i\subseteq\mathcal S_n\) and a selector
\(\iota:\mathcal S_n\to[m]\), and use the refined report of
Section~\ref{sec:reporter}. In particular, an all-pass report reveals
all leaf sums, and an ordinary report reveals the central sum and all
unselected leaf sums. The total number of input bits is
\(N=n(m+1)\).
\end{definition}

The tests are unions of complete sum levels. They cannot distinguish
two sign strings having the same sum. This restriction is part of the
definition and part of the minimality conclusion.

\begin{lemma}[Exact degree of a nonconstant reporter]
\label{enum:exact-degree}
Every nonconstant reporter in Definition~\ref{enum:definition} has
maximum cell degree exactly \(nm\).
\end{lemma}

\begin{proof}
Proposition~\ref{prop:reporter} gives the upper bound. The indicator
of any nonempty \(n\)-bit sum level has degree exactly \(n\): its
top Walsh character is constant, with value
\((-1)^{(n-x_j)/2}\), throughout that level, so the corresponding
coefficient is that sign times \(p_j\), and is nonzero.

If all tests are nonempty, choose one allowed value in each test.
The corresponding all-pass cell fixes these \(m\) raw leaf sums and
places no condition on the center. Its indicator is a product of
\(m\) level indicators on disjoint blocks. Their nonzero top
coefficients multiply, giving degree \(nm\).

If a test is empty, the all-pass event is empty. Unless the whole
observation is the single exceptional tag, there is an attained
ordinary cell. Such a cell fixes the central sum and \(m-1\)
unselected leaf sums; its reported values already force an unselected
failure. Its indicator is again a product of \(m\) level indicators,
and has degree \(nm\).
\end{proof}

Thus an advantage over the block bound \(nm\) is an advantage over
the actual maximum cell degree, apart from constant observations,
which have no advantage.

\subsection{The exact finite objective}

The singleton formula of Theorem~\ref{thm:finite-selector} extends
to arbitrary nonempty proper tests without numerical integration.
For test \(I_i\), put
\[
 \beta_i=\Pbb(Z_i\in I_i),\quad
 \mu_i=\E[Z_i\mid Z_i\in I_i],\quad
 s_i^2=\Var(Z_i\mid Z_i\in I_i).
\]
For a selector \(j\mapsto i_j=\iota(x_j)\), define
\begin{align*}
 Q&=\prod_i\beta_i,& S&=\sum_i s_i^2,&
 w_j&=p_j/\beta_{i_j},\\
 R_0&=\sum_jw_j,& K&=R_0-1,&
 R_1&=\sum_jw_j(x_j-\mu_{i_j}),\\
 H_2&=\sum_jw_j(x_j-\mu_{i_j})^2.&&&&
\end{align*}
Every number here is rational. Since all tests are proper,
\(\beta_i<1\), and hence \(K>0\).

\begin{proposition}[Objective for arbitrary sum-level tests]
\label{enum:objective}
For the reporter \(F\), its exact advantage is
\begin{equation}\label{enum:gain}
 g:=v(F)-nm
 =Q\left\{\frac{R_1^2}{K}-H_2-KS+
               \sum_j w_j s_{i_j}^2\right\}.
\end{equation}
The exceptional cell has probability \(QK\) and optimal predictor
\[
 b_* = \sum_i\mu_i+\frac{R_1}{K}.
\]
\end{proposition}

\begin{proof}
The identities \(\Pbb(P)=QR_0\), \(\Pbb(B)=Q\), and
\(\Pbb(A)=QK\) follow by conditioning on the center.
For a constant exceptional predictor \(b\), put
\(C=\sum_i\mu_i-b\). Conditional on \(U=x_j\) and \(P\), the
mean of \(L-b\) is \(x_j+C-\mu_{i_j}\), and its variance is
\(n+S-s_{i_j}^2\). On \(B\), its mean is \(C\) and its variance
is \(n+S\). Subtracting the contribution of \(B\) from that of
\(P\) gives
\[
 \E[\1_A\{n-(L-b)^2\}]
 =Q\left\{-KC^2-2R_1C-H_2-KS+
                       \sum_jw_js_{i_j}^2\right\}.
\]
The maximum occurs at \(C=-R_1/K\). The optimizing predictor is
the conditional mean on \(A\), so
\eqref{eq:optimal-exceptional} identifies the maximum with
\(v(F)-nm\). This proves \eqref{enum:gain}.
\end{proof}

\subsection{Reductions that make the enumeration complete}

\begin{lemma}[Reduction to distinct proper tests and onto selectors]
\label{enum:reductions}
If a reporter has positive advantage, then it has no empty or full
leaf test. From any positive reporter one can obtain another positive
reporter with no more input bits, distinct tests, and a surjective
selector. A reporter with one leaf has no positive advantage.
\end{lemma}

\begin{proof}
If at least two tests are empty, \(A\) is empty. If exactly one test,
at leaf \(i\), is empty, then \(A\) requires selecting \(i\) and
the success of all other tests. Thus \(A\) is independent of
\(Z_i\), which retains variance \(n\) on that event. In either
case \eqref{eq:optimal-exceptional} rules out a positive advantage.

If test \(i\) is full, then \(A\) is false when \(i\) is selected,
and is independent of \(Z_i\) when another leaf is selected.
Consequently \(A\) is independent of \(Z_i\) altogether. Again
\(\Var(L\mid A)\ge n\), whenever \(A\) is nonempty. This also
rules out using a full test as a favorable selector fallback.

Suppose next that a proper nonempty leaf test is never selected.
Deleting it leaves \(K,R_1,H_2\) and the selected-variance term in
\eqref{enum:gain} unchanged, decreases \(S\) by \(s_i^2\), and
divides \(Q\) by \(\beta_i\). The new advantage is
\[
 g'=\frac{g}{\beta_i}+\frac{QK}{\beta_i}s_i^2>0.
\]
If two tests coincide, merge their selector labels. The parameters
in \eqref{enum:gain} do not change, because the two tests have the
same probability, mean, and variance. One copy is now unused and
can be deleted. Repeating these operations gives the stated
reduction. A surjective selector has at most \(n+1\) leaf labels.

Finally, when \(m=1\), the reporter only describes its leaf and
never the independent center. Its retained variance is at most
\(n\). By Lemma~\ref{enum:exact-degree}, a nonconstant such
reporter has degree \(n\); a constant one has both variance and
degree zero.
\end{proof}

\subsection{The finite certificate}

\begin{theorem}[Twenty-bit minimality in the binomial reporter class]
\label{enum:minimality}
No reporter in Definition~\ref{enum:definition} on fewer than twenty
bits has retained variance greater than its maximum cell degree.
The reporter of Theorem~\ref{thm:witness} attains a strict advantage
on twenty bits. Thus twenty is the minimum within this specified
class.
\end{theorem}

\begin{proof}
It suffices to consider a nonconstant reporter with distinct nonempty
proper tests and a surjective selector, by
Lemmas~\ref{enum:exact-degree} and~\ref{enum:reductions}.
Its codimension is exactly \(n\). If \(n\le3\), the analytic
codimension obstruction excludes an advantage. The case \(m=1\)
has already been excluded. If \(n\ge7\) and \(m\ge2\), then
\(n(m+1)\ge21\). Therefore fewer than twenty bits leave only the
four parameter pairs in the following table.

\begin{center}
\begin{tabular}{@{}rrrrr@{}}
\toprule
Block size \(n\)&Leaves \(m\)&Bits \(N\)&Configurations&Maximum \(g\)\\
\midrule
4&2&12&13,050&\(-1/12\)\\
4&3&16&609,000&\(-35/17152\)\\
5&2&15&117,242&\(-595/8576\)\\
6&2&18&992,250&\(-1293/33152\)\\
\bottomrule
\end{tabular}
\end{center}

There are \(q_n=2^{n+1}-2\) nonempty proper sum-level tests. The
certificate enumerates every unordered family of \(m\) distinct
tests and every onto map from the \(n+1\) central sums to those
tests. The exact number of configurations is therefore
\begin{equation}\label{enum:count}
 \binom{q_n}{m}
 \sum_{r=0}^m(-1)^r\binom mr(m-r)^{n+1}.
\end{equation}
This gives exactly the counts in the table. Sorting the test masks
only removes permutations of leaf labels; all selectors for each
sorted family are still included.

The program \texttt{code/enumerate\_reporters.py} evaluates
\eqref{enum:gain} with integer fractions for every configuration,
compares them exactly, and obtains the displayed maxima. All
1,731,542 configurations have nonpositive advantage; in these four
reduced cases the maxima are strictly negative. The JSON output
records the counts and a maximizing family, selector, and predictor
for each row. This exact finite exhaustion completes the lower
bound. The twenty-bit upper bound is
Theorem~\ref{thm:witness}.
\end{proof}

The only computational part of this theorem is the four-row finite
exhaustion. The codimension threshold and the positive twenty-bit
construction have independent mathematical proofs. In particular,
their validity does not depend on this minimality search.

\subsection{Reproduction and limits of the certificate}

The companion code uses only the Python standard library. The command
\begin{verbatim}
python code/run_verification.py
\end{verbatim}
checks the twenty-bit witness and writes
\texttt{data/witness\_certificate.json}. It enumerates all \(5^5=3125\)
tuples of block sums, with their exact binomial multiplicities, so
every one of the \(2^{20}\) sign strings is accounted for. All 622
cells, their conditional moments, the complete score law, and the
retained variance are checked.

The same verifier also performs a complete degree check. Block symmetry
makes a Walsh coefficient depend only on the five subset sizes in its
blocks. There are 56 such profiles of total degree greater than 16;
all their coefficients vanish on each of the 622 cells, giving 34,832
exact zero checks. A degree-16 coefficient is nonzero on every cell.
The code uses the exact sum of a fixed character on each binomial
layer, rather than a numerical transform.

To regenerate the complete minimality certificate, run
\begin{verbatim}
python code/run_verification.py --exhaustive
\end{verbatim}
This writes \nolinkurl{data/minimality_certificate.json}. A separate
implementation, \nolinkurl{code/enumerate_two_leaf_integer.py},
uses cross-multiplied integer comparisons in its inner loop. The
included independent check of the largest two-leaf case agrees with
the fraction implementation on all 992,250 configurations and on
the exact maximum \(-1293/33152\).

These are reproducible finite certificates, not a Lean formalization
of the enumeration programs. They do not exclude smaller witnesses
outside Definition~\ref{enum:definition}, including leaf tests that
split individual sum levels, unequal block sizes, different reporting
mechanisms, or general partitions of a cube.
